\section{Introduction}
\label{sec:introduction}

An autoregressive (AR) language model such as GPT-2 predicts the next token from the text so far. A masked diffusion language model (DLM) starts with blank positions and can fill several of them in a pass. We call each pass a denoising step. \citet{gong2025diffullama} trained existing AR models such as GPT-2 and LLaMA-2 to fill in these blanks. We call the original model the AR base. The authors hypothesise that ``the adapted model retains some of the abilities from the base AR model''.

Several attempts at the same question can contain a correct solution even when the first one is wrong. On SATMath, DiffuLLaMA answers 23.6\% correctly when the prompt includes worked examples. Majority vote over three samples raises accuracy to 27.7\%, a method known as self-consistency \citep{wang2023selfconsistency}. At least one sample is correct 57.7\% of the time. The authors attribute this gap to ``high uncertainty about its responses''. How to choose better answers remains open.

We test two ways to choose from the same set of solutions, or candidates. Majority vote picks the final number that appears most often, even when the solutions reaching it use different reasoning. In AR reranking, an AR model scores each solution with its next-token probabilities and picks the highest-scoring one. The model's weights stay fixed. We vary the number of denoising steps to see how the choice of answer changes as the generated solutions improve. We ask whether reranking gains are larger when the DLM and its AR base have similar hidden activations on the same text. We call that similarity representation preservation and measure it with centred kernel alignment (CKA) \citep{kornblith2019cka}. Section~\ref{sec:formulation} defines the score.

\subsection{Contributions}

\lead{A shared set of solutions}
We evaluate DiffuGPT-S, DiffuGPT-M, DiffuLLaMA \citep{gong2025diffullama}, and Dream 7B \citep{ye2025dream7b}, each adapted from an AR base. LLaDA-8B, which was trained from scratch, is the control \citep{nie2025llada}. We save $N = 10$ solutions per model for each GSM8K math problem \citep{cobbe2021gsm8k}. All selection rules see the same solutions. They have no access to the reference answers. The first solution is the baseline for gains. Oracle accuracy is the fraction of questions with at least one correct solution. It gives the best possible score on those solutions.

\lead{Answer selection across generation settings}
Going from 16 to 32 steps improves the first solution more than voting at 16 steps. From 16 to 64 steps, voting closes a growing share of the gap to oracle accuracy. On matched questions from 64 to 256 steps, that share falls from 69\% to 53\% on Dream and from 68\% to 40\% on LLaDA. Voting still improves accuracy at 256 steps. If the choice is between one 256-step greedy solution and voting over ten 64-step solutions, Dream is 9.0 points more accurate with one solution and LLaDA is 3.5 points more accurate with voting.

\lead{Scoring and representation similarity}
At 64 steps on 500 questions, AR reranking beats the first solution on Dream and LLaDA and trails voting by 23.6 and 21.0 points, respectively. At 256 steps it trails by 8.6 and 2.6 points on the matched subsets. LLaDA's confidence interval includes zero. Scoring through the final answer raises Dream's reranking accuracy by 13 points with Qwen2.5-7B. It also changes the apparent relationship between model similarity and reranking gains (Section~\ref{sec:experiments}).
